
\chapter{引言}

对于\glsfirst{ml}爱好者来说，这是一个令人振奋的时代。
\glspl{ml}——构建学习机器的科学——正以前所未有的速度催生新算法。
\gls{ml}也在日常生活中占据更大比重，即使远离该领域的人也很难不注意到这一点。

为说明这一发展的范围和速度，分享以下个人经历。
在开始撰写本论文前不久，我了解到新的 DALL-E \cite{ramesh2022hierarchical} 算法，它能够根据用户提示生成高分辨率图像。
这让我震惊于\gls{ml}社区的能力，Twitter 上的大多数人也是如此，他们纷纷尝试创作最佳提示。
除了精彩的模因和笑话之外，这类技术的前景和潜在应用——以及潜在风险——超乎我的想象。
仅几周后，我便得知不同研究团队已开发出类似算法，如 Stable Diffusion \cite{rombach2021highresolution} 和 Imagen \cite{saharia2022photorealistic}。
我曾以为自己有幸见证了\gls{ml}这一新方向的诞生。
这一信念直到前几天还在——在撰写本论文时，我偶然发现了 Make-a-Video \cite{singer2022make}。
该算法可以根据用户提示生成高分辨率视频……
仅几天后，Imagen 的视频扩展版——Imagen Video \cite{ho2022imagen} 也发布了。
我还没来得及思考图像生成算法的全部潜力，视频生成算法就已诞生。
我无法想象这个领域的下一个发展将是什么\footnote{这还没提到 ChatGPT，它是在我修改论文时发布的。}。

本论文的核心主题——强化学习（Reinforcement Learning, RL）\cite{sutton2018reinforcement}——是\gls{ml}的一个子领域，知名度较低但同样有其显著成就。
现在让我们深入这个充满活力的子领域。

\section{强化学习范式}

\gls{rl}的主要目标是构想一个\emph{智能体（agent）}，通过与其\emph{环境（environment）}交互学习最大化某种效用度量。交互过程建模为序贯决策问题。每一步，智能体在环境中执行一个\emph{动作（action）}。作为该动作和外部因素的结果，环境转换到新\emph{状态（state）}并向智能体提供强化信号。该信号称为\emph{奖励（reward）}，对动作质量提供反馈。智能体的目标是最大化累积奖励和——即\emph{回报（return）}。环境遵循的序贯过程（包括转移动态和奖励）由\glsfirst{mdp}建模。为最大化奖励，智能体必须自主平衡探索（exploration）和利用（exploitation）。它需探索\gls{mdp}的新区域以防止错过机会，同时考虑利用已知有利区域以最大化效用。

\gls{rl}框架受心理学关于人类学习过程发现的启发。它足够通用以涵盖许多问题，特别是\gls{ml}的其他子领域——\emph{监督学习（supervised learning）}和\emph{无监督学习（unsupervised learning）}。它也足够具体以产生有意义的理论结果指导算法设计。

在实证方面，框架的有效性由其带来的突破性进展所支持。其中部分成就包括：
机器人运动\cite{haarnoja2018learning}；
自动驾驶\cite{kiran2021deep}；
在专家级别玩复杂视频游戏\cite{vinyals2019grandmaster}；
设计性能更高且能耗更低的计算机芯片\cite{mirhoseini2021graph}；
控制托卡马克中等离子体形状\cite{degrave2022magnetic}。

然而，大多数这些应用需要对原始\gls{rl}框架进行轻微扩展，例如考虑状态的部分可观测性。原始框架的另一个限制可能是动力系统中存在延迟。下一节将解释为何考虑延迟是实现出色实证结果的关键。尽管\gls{ml}和\gls{rl}研究进展如此迅速，但我们建议读者花些时间深入探讨延迟问题。

\section{强化学习中的延迟}
\label{sec:why_delay_ns}

如前所述，\gls{rl}及其\gls{mdp}模型是序贯决策的优质框架。然而，如读者所知，模型是对真实过程的简化，有其自身局限。审视\gls{rl}实际应用的文献可揭示这些局限所在。\cite{mahmood2018setting} 发现，在真实机器人应用中，从一系列候选障碍中，学习的两大主要障碍是动作空间的选择和状态观测或动作执行中的延迟。这些延迟通过向智能体提供过时信息或延迟智能体对环境的控制来影响智能体-环境交互。这是支持延迟研究的有力论据。

我们认为延迟是任何序贯问题固有的。人类面临的最普遍延迟类型是大脑处理来自身体传感器信息所需的时间。例如，运动信息从光感受器传递到大脑高级视觉区域所需时间估计约为 100 ms \cite{de1991vernier}。在\gls{rl}中，此延迟类似于传感器采集时间导致的延迟。延迟也可能来自智能体根据环境观测计算下一个动作所需的时间。例如，使用更复杂的策略（如\emph{深度神经网络（deep neural network）}）会加剧延迟，这在主动流动控制中已有观察到\cite{ren2021applying,mao2022active}。

延迟自然出现的另一个例子是交易。在实践中，信息收集和交易传输的速度极其重要。信息通过光纤走不同路径，导致在不同地点的异步到达时间。这些延迟可被利用以获利\cite{lewis2014flash}。此外，通过抽象掉市场上的其他智能体，单个交易智能体的回报受到延迟的负面影响\cite{wilcox1993effect}。使用\gls{rl}智能体的外汇（FX）实验证实了引入延迟后的性能下降\cite{liotet2022delayed}。延迟影响的其他例子包括带半挂车的拖拉机紧急制动时间\cite{bayan2010brake}；并行计算中计算节点可用信息\cite{hannah2018unbounded}；机械臂控制\cite{mahmood2018setting}。

即使环境看似同步运行（如围棋），延迟仍隐藏在某处。在关于 AlphaGo \cite{silver2016mastering}（与世界冠军李世石对弈的算法）开发历程的纪录片中，AlphaGo 花费了超过 30 秒计算第一步。围棋比赛总时间有限，过长的动作计算可能有害。在此例中延迟并未造成损害，因为 AlphaGo 仍有充足时间并最终战胜了李世石。

在仿真中，延迟可以人为去除，但这可能带走重要的真实感。在 Atari 2600 等游戏中，延迟显著影响\gls{rl}智能体的性能\cite{firoiu2018human}。一个简单解决方案是在智能体选择动作时暂停环境以消除延迟。然而，\cite{firoiu2018human} 认为，由于大多数游戏以 60Hz 显示图像，且研究发现某些人群的平均反应时间约为 250 ms \cite{jain2015comparative}，不为人工智能体建模延迟使其相对于人类获得不公平优势。

许多相关领域也考虑延迟问题。例如，在\emph{控制理论（control theory）}中\cite{chung1995time, dugard1998stability, niculescu2001delay, gu2003survey, richard2003time, fridman2014introduction}，延迟对系统稳定性有复杂影响\cite{kolmanovskii1999delay}。该领域特别考虑了拥塞控制问题，其中数据包从多个源发送到网络\cite{jacobson1988congestion}。在此问题中，延迟自然产生，因为数据包从不同源走不同路径\cite{witsenhausen1971separation, altman1999congestion,ait2003stability}。在\emph{在线学习（online learning）}中，接收延迟反馈会减慢学习过程\cite{joulani2013online,pike2018bandits,lancewicki2021stochastic}。在\emph{实时（real-time）}\gls{rl}中，智能体和环境之间的交互不再是同步的——即环境"等待"智能体选择动作——而是变成异步的。实际上，智能体必须在环境持续演变的同时计算下一个动作，且无法保证计算出的动作对新的环境状态仍然相关。有趣的是，该子领域已被证明等价于延迟\gls{rl}\cite[Theorem~1]{ramstedt2019real}。因此，该领域的兴趣\cite{hester2013texplore,caarls2015parallel,travnik2018reactive,ramstedt2019real,xiao2020thinking}也论证了延迟研究的价值。

基于所有上述原因，我们认为延迟是\gls{rl}实际应用中的核心研究方向。此外，在传统框架中引入延迟所引发的理论问题也具有研究价值，因为它们与\gls{rl}的许多子领域相关，如\glsfirst{pomdp}或实时\gls{rl}。

\section{原创贡献}

本论文聚焦\gls{rl}场景中特有的两类延迟：状态观测延迟（state observation delay）和动作执行延迟（action execution delay）。目标是拓宽对此类延迟的理解并提供高效解决方案。特别关注文献中普遍存在且在许多情况下能很好近似实际延迟的常值延迟（constant delay）。然而，我们将反复研究将框架扩展到更原始延迟类型的理论和实证影响。本论文的贡献遵循四个目标：提供\gls{rl}中延迟的统一框架；为更好理解延迟提供理论基础；设计有理论依据且高效的算法以应对延迟环境；对所提算法及文献中算法进行广泛实证评估。下文将详述每个目标的贡献，然后引用本人合著的相关出版物。

\subsubsection{延迟的统一框架}

本文对\gls{rl}及相关领域的延迟文献进行了广泛综述。重点放在\gls{rl}中的状态观测和动作执行延迟，但也讨论了其他类型延迟。值得注意的是，我们将文献中的算法归纳为三类。第一类是\emph{增广状态（augmented state）}方法。其思想是将最后观测到的状态与智能体已选择但尚未观察到结果的动作序列进行增广。使用这种新状态可将延迟问题重新转化为传统的无延迟\gls{mdp}框架。第二类是\emph{无记忆（memoryless）}方法，智能体忽略延迟，仅使用最后观测状态的信息。这一思想简单，但在实践中可能高效。最后，\emph{基于模型（model-based）}方法学习环境模型以预测动作将施加时的环境状态。据我们所知，此前尚无工作提供\gls{rl}中延迟的综述。我们认为提供此信息对该领域有益，因为文献中经常出现重叠结果。

此外，我们提出延迟问题的统一数学框架。这一提议遵循同一思路，即延迟文献可受益于更多结构化。我们的框架将所有最常見的延迟类型作为特例包含在内，并与\gls{mdp}模型自然衔接。

\subsubsection{延迟的理论理解}

首先，分析延迟状态观测或动作执行对\gls{mdp}的理论影响。第一个结果是正式证明一个看似显而易见但尚未被证明的事实：延迟越长，最大可达累积奖励（即回报）越低。同理，展示延迟如何影响策略集所能实现的回报变异性。更长的延迟缩小了智能体可获得回报的范围，但也降低了重复学习同一任务时回报的方差。然后，将常值延迟框架扩展到多动作延迟（multiple action delay）框架，其中单个动作可影响未来的多个转移。该框架从理论角度具有研究价值，因为它将常延迟过程作为特例包含在内。此外，它自然地将高阶\glspl{mc}模型扩展到\glspl{mdp}。我们展示多动作延迟过程可通过状态增广重新转化为\gls{mdp}，与常延迟过程类似。然后研究该框架相对于传统\gls{mdp}的特性和特殊性。特别是，这有助于强调高阶\gls{mdp}与延迟\gls{mdp}之间的联系。值得注意的是，一些性质（如单链性（unichain））不会从底层\gls{mdp}传递到基于它构建的延迟过程，而其他性质（如连通性（communicating））则会传递。

其次，将焦点转向延迟\gls{rl}算法并研究其保证。分析处理延迟的三种方法：增广状态、无记忆和基于模型的方法。已知增广状态方法可为延迟问题找到最优策略。相反，我们展示基于模型方法可能过于严格而丢弃最优策略。然而，在平滑\gls{mdp}（smooth MDP）——将在本文中定义的概念——背景下，我们证明基于模型方法能提供接近最优的策略。将上述保证扩展到非整数延迟（non-integer delay）。非整数延迟让人联想到实时\gls{rl}的异步环境假设。智能体以给定频率观测状态，但只能以相位变化的方式执行动作。最后，推导出基于模型策略在常值延迟上训练但在随机延迟上测试时的性能界。

\subsubsection{延迟强化学习算法}

在算法方面，首先探索基于模型方法，设计一种能够学习当前状态概率分布（即信念（belief））的向量化表示的算法。为学习信念表示，第一个神经网络（Transformer）处理增广状态并为每个未来状态生成表示（直至无延迟状态）。该表示通过拟合分布的另一个网络（归一化流（normalising flow）网络）训练以对应信念。将该表示接入\gls{rl}算法可提升性能并加速学习速度，效果在随机环境中尤为明显。此外，提供一种直接方法，使预训练于特定延迟的该算法能适应更短延迟。

然后探索一种不同的方法，通过模仿无延迟专家行为来学习策略。实际中，延迟智能体被训练为复制无延迟智能体基于当前状态选择的动作，但仅使用延迟信息——即增广状态。展示如何利用平滑\glspl{mdp}的理论保证与模仿学习（imitation learning）的理论结果相结合来指导算法设计。实证表明该方法达到最先进性能，且足够通用以适用于常值、非整数和随机延迟。

作为第三种方法，探索无记忆方法。在此情形下，问题本质上变成非平稳（non-stationary）。因此，设计一种理论，通过优化未来性能估计的概率下界来适应未来非平稳性。该未来性能估计基于多重重要性采样（importance sampling）以考虑非平稳性导致的过去动态多样性。在平滑非平稳条件下，该估计器的偏差有界。为进一步考虑过程的随机性，利用浓度不等式获得估计的下界。这涉及前一估计器的方差上界。有趣的是，控制估计器的方差间接约束了策略的非平稳性。这防止对过去非平稳动态的过拟合并更好地泛化到未来动态。实际中，算法实现为一个超策略（hyper-policy），接受时间作为输入并输出每个时刻要查询的策略。实验展示该算法如何学习过程的非平稳性以适应它，包括当非平稳性由延迟引起时。

\subsubsection{广泛实证评估}

最后一个贡献是对本文提出的算法及文献中算法进行广泛实证评估。设计多种测试任务，从平衡摆到货币汇率交易。在这些任务中，改变延迟的长度和性质以研究算法的鲁棒性。考虑常值、非整数、随机和多动作延迟。这些实证评估有助于理解算法的特性和偏好设置。

\subsubsection{相关出版物}

上述贡献可在以下论文中找到。\cite{liotet2021learning} 中提出基于模型的方法并分析其可能过于严格的问题。在该工作中，我设计了方法，进行了理论分析和大部分实验分析。模仿学习方法及平滑\glspl{mdp}中延迟的分析见 \cite{liotet2022delayed}。我完成了常值延迟的部分理论分析以及向非整数和随机延迟的扩展，并进行了大部分实验分析。非平稳无记忆方法来自 \cite{liotet2022lifelong}。在该工作中，我进行了大部分理论分析、部分方法设计以及全部实证研究。最后，关于多动作延迟扩展及改变延迟分布对智能体性能影响的分析结果将作为待提交论文的基础。在该工作中，我完成了全部理论和实证分析。

\section{论文概述}

本论文前两章专注于介绍和定义\gls{rl}的整体概念以及延迟文献中的特定概念。

\begin{itemize}
    \item \Cref{chap:prelim} 介绍后续章节使用的主要数学工具和算法素材。特别地，介绍\gls{mdp}模型和\gls{rl}的一般框架。
    \item \Cref{chap:delay} 深入探讨作为第一章所述原始\gls{rl}框架扩展的延迟问题。引入相关符号和术语。随后详尽列举文献和实际应用中遇到的延迟类型。最后讨论\gls{rl}及相关领域的延迟文献。
\end{itemize}

\noindent 这些介绍性章节之后是四个章节，展示本论文的原创贡献。每章结构相同：第一部分正式介绍问题并详述提出的解决方案；第二部分提供理论分析以更好理解问题特性和解决方案的性质；第三部分通过彻底的实证分析支持前两部分的论断。

\begin{itemize}
    \item \Cref{chap:belief_based} 聚焦状态观测和动作执行中的常值延迟。提出一种称为\emph{信念表征网络（Belief Representation Network）}的概率方法，其中智能体预测近期未来以应对延迟。理论分析将提供对这类延迟问题及其对性能负面影响的理解。本章内容见 \cite{liotet2021learning}。
    \item \Cref{chap:imitation_undelayed} 与前一章处于相同框架。提出一个简单解决方案——{DIDA}——延迟智能体学习模仿无延迟专家行为。尽管简单，但证明了 DIDA 在平滑环境中的强大理论保证。该解决方案扩展到新框架，包括非整数和随机延迟。本章内容源自 \cite{liotet2022delayed}。
    \item \Cref{chap:lifelong} 同样设定在常值延迟框架中，提出另一种方法。解决方案考虑无记忆智能体，即对延迟不可见。从智能体角度看，过程变为非平稳，因此设计非平稳策略以适应演化动态。通过优化智能体未来性能的估计来实现，称该方法为 POLIS。本章通过研究未来性能估计器的偏差和方差为 POLIS 设计提供理论基础。本章内容取自 \cite{liotet2022lifelong}。
    \item \Cref{chap:multi_action_delay} 考虑包含常值延迟的延迟框架。其中，一个动作的执行可延伸至多个未来步骤，从而产生多动作延迟（multiple action delay）。深入分析新框架的性质，特别关注理解延迟分布对智能体最佳性能的影响。本章材料尚未在其他地方发表。
\end{itemize}

最后，\Cref{chap:conclusion} 总结论文的主要成果和要点以及局限性。重要的是，每个结果都将引导读者到论文中呈现该结果的部分。本章还将讨论潜在研究方向。

更多证明、实验细节和实证结果见附录 \Cref{app:proofs_results}。

